\scanpage{231}
Par Eilenberg-Zilber, la \deleted{co}homologie de
$C_*(A_*\otimes_k B_*)$ (\deleted{deux}\inserted{$A_*,B_*$}
$k$-modules $1/2$ simpliciaux \deleted{$\unsure{\Phi_*\longrightarrow}$}
\deleted{complexe de chaines}\inserted{op. bord}
$\partial=(\deleted{\unsure{\sum}}\partial_n)$,
$\partial_n=\sum_0^n\partial^*_{i,n}$,
\margin{$C_n(A_*\otimes B_*)=A_n\otimes B_n$}
c'est donc, par Dold-Puppe, aussi $\pi_*(A_*\otimes B_*)$)
se calcule en prenant \deleted{du}\inserted{l'homologie} produit tensoriel
\deleted{des complexes} $C_*(A_*)\otimes C_*(B_*)$, qui est lui-même
\inserted{``chaine-}\unsure{homotope|homotopie}'', par Dold-Puppe, à
$\deleted{\unsure{C_*(A}}A_*^!\otimes B_*^!$
(car $A_*^!\xrightarrow[\text{chaine-homotopisme}]{\sim}A_*$,
$B_*^!\mathbin{-}B_*$);
\deleted{lequel}\inserted{dont l'homologie} peut se calculer par Künneth.
Ici $B_*^!=(\Lambda^i\Psi_*)^!\simeq k[i]$, donc on \deleted{\unsure{dit on}} a
\[
\underset{\substack{\Vert\\\pi_\ell(A_*\otimes B_*)}}{H_\ell(A_*^!\otimes B_*^!)}
\simeq
\underset{\substack{\Vert\\\pi_{\ell-i}(A_*\otimes B_*)}}{H_{\ell-i}(A_*^!)}.
\]

\uline{Cor.} 1.5
\margin{\figure{\begin{tikzpicture}\draw (0,0)--(0,1.4);\end{tikzpicture}}}
Considérons p.ex. le complexe \deleted{de chaines}\inserted{$1/2$ simplicial}
De Rham à puissances divisées
\[
\mathcal A_*^{\overset{\inserted{\text{degré total}}}{*}\,\overset{\inserted{\text{degré ext}}}{*}}
\simeq\Gamma^*(\Phi_*)\otimes_k\Lambda^*\Psi_*
\]
on trouve alors
\[
\pi_\ell(\mathcal A_*^{*i})\simeq
\pi_{\ell-i}(\underbrace{\mathcal A_*^{*0}}_{\Gamma^*(\Phi_*)})
\xleftarrow{\sim}\pi_{\ell-i}(\underbrace{\mathcal A_*^{00}}_{\Gamma^*(0)=k_*})
=\begin{cases}0&\text{si }\ell\ne i\\k&\text{si }\ell=i.\end{cases}
\]
En \deleted{\unsure{tous}} particulier
\[
\pi_\ell(\underbrace{\mathcal A_*^{p,i}}_{\Gamma_k^{p-i}\Phi_*\otimes\Lambda^i\Psi_*})
\simeq\pi_{\ell-i}(\deleted{\unsure{\Gamma^i}}
\underbrace{\mathcal A_*^{\unsure{p-i,0|p-i0}}}_{\Gamma_*^{p-i}\Phi_*})
\simeq\begin{cases}0&\text{si }\ell\ne i\text{ ou }p\ne i\\k&\text{si }i=\ell=p\end{cases}
\]
donc, pour $p$ fixé, $i$ variable, on trouve que les \inserted{composantes}
$\mathcal A_*^{p,0},\mathcal A_*^{p,1}\mathbin{-}\mathcal A_*^{p,p-1}$
\deleted{de} (en degré $<p$) de $\mathcal A_*^{p*}$ \deleted{sont}
(qui est une résolution de $k_*$) sont flasques, la seule composante non flasque étant
$\mathcal A_*^{p,p}$ (ayant comme seul gr. d'hom. non nul
$\pi_p(\mathcal A_*^{p,p})=k$). On peut donc considérer que
\margin{230}
